\documentclass[11pt]{article}
\usepackage[a4paper,margin=2.3cm]{geometry}
\usepackage{amsmath,amssymb,bm,mathtools}
\usepackage{enumitem}
\usepackage[dvipsnames]{xcolor}
\usepackage[most]{tcolorbox}
\usepackage{microtype}
\usepackage[colorlinks=true,linkcolor=NavyBlue,citecolor=ForestGreen,urlcolor=NavyBlue]{hyperref}
\newtcolorbox{keybox}[1][]{colback=Gray!6,colframe=Gray!55,boxrule=0.5pt,arc=2pt,left=6pt,right=6pt,top=5pt,bottom=5pt,#1}
\newtcolorbox{propbox}[1][]{colback=NavyBlue!5,colframe=NavyBlue!55,boxrule=0.6pt,arc=2pt,left=7pt,right=7pt,top=6pt,bottom=6pt,#1}
\newtcolorbox{algbox}[1][]{colback=ForestGreen!5,colframe=ForestGreen!50,boxrule=0.6pt,arc=2pt,left=7pt,right=7pt,top=6pt,bottom=6pt,#1}
\newcommand{\herm}{^{\mathsf{H}}}
\newcommand{\vh}{\mathbf{h}}\newcommand{\vv}{\mathbf{v}}\newcommand{\mH}{\mathbf{H}}
\newcommand{\mC}{\mathbf{C}}\newcommand{\mI}{\mathbf{I}}\newcommand{\mR}{\mathbf{R}}\newcommand{\mQ}{\mathbf{Q}}
\newcommand{\Sk}{\mathcal{S}_k}\newcommand{\E}{\mathbb{E}}
\title{\textbf{A Detector-Aware Information-Surplus Clustering Rule (P1)\\for Cell-Free Massive MIMO: Principle, Steps, and Analysis}}
\date{\today}
\begin{document}\maketitle
\newcommand{\vexp}[2]{\par\smallskip\noindent\textit{#2}}

\begin{keybox}
\textbf{In one line.} P1 chooses, for each user, the \emph{smallest} set of access points (APs)
whose post-combining SINRs add up to at least a fraction $(1-\epsilon)$ of the user's total SINR
over all APs. This gives a \emph{tunable, provable} guarantee on how much information the cluster
keeps, using the quantity the detector actually sees (SINR), not just the channel strength.
\end{keybox}

\section{Background and motivation (why cluster at all)}
In cell-free massive MIMO a user is served not by all $L$ APs but by a small \emph{user-centric}
cluster $\Sk\subseteq\{1,\dots,L\}$, to keep the fronthaul and the receiver complexity scalable
\cite{ngo2017cellfree,buzzi2017usercentric,interdonato2019ubiquitous,bjornson2020scalable}. The open
question is \emph{which} APs to pick. Two common families are:
\begin{itemize}[leftmargin=1.4em,itemsep=1pt]
\item \textbf{Channel-norm / large-scale-fading (CN):} keep the APs whose channel strength exceeds a
      threshold (dynamic cooperation clustering, largest-large-scale-fading selection)
      \cite{bjornson2020scalable,ngo2018totalenergy}. Simple, but ignores interference and
      estimation error.
\item \textbf{Information-rate (IR):} keep the APs whose per-link achievable rate exceeds a threshold
      \cite{mashdour2024clustering}. Better, but the rate used omits multi-user interference, and the
      cluster size grows with SNR (hence the fronthaul grows).
\end{itemize}
\textbf{Our aim.} A rule that (i) uses the \emph{receiver} quantity (post-combining SINR, which
already contains interference and estimation error), (ii) gives a \emph{provable} bound on the
information lost by clustering, and (iii) keeps a \emph{bounded, stable} cluster size. This is P1.

\section{The quantity P1 acts on: per-AP SINR}
Each AP $l$ forms a local MMSE combiner for user $k$ from the estimated channels
\cite{nayebi2016lsfd,bjornson2020scalable}
\begin{equation}
\vv_{l,k}=\big(p\,\hat{\mH}_{l}\hat{\mH}_{l}\herm+\mQ_{l}\big)^{-1}\hat{\vh}_{l,k},\qquad
\mQ_{l}=p\sum_{j=1}^{K}\mC_{l,j}+\mI_{N},
\label{eq:comb}
\end{equation}
\vexp{$\vv_{l,k}$ is the AP-$l$ receive filter for user $k$; $\hat{\mH}_l=[\hat\vh_{l,1},\dots,\hat\vh_{l,K}]$ the estimates; $p$ the SNR; $\mC_{l,j}$ the estimation-error covariance; $\mQ_l$ the noise-plus-error term; $N$ antennas per AP.}%
{Here $\vv_{l,k}$ is the receive filter AP $l$ uses for user $k$. $\hat{\mH}_{l}=[\hat{\vh}_{l,1},
\dots,\hat{\vh}_{l,K}]$ stacks the estimated channels of all $K$ users at AP $l$, $p$ is the SNR,
$\mC_{l,j}$ is the channel-estimation-error covariance of user $j$, and $\mQ_{l}$ collects the noise
($\mI_{N}$) and the error energy. $N$ is the number of antennas per AP. In words: each AP builds the
best linear filter it can from what it \emph{estimated}, not from the perfect channel.}
The resulting post-combining SINR of user $k$ at AP $l$ is
\begin{equation}
\gamma_{l,k}=\frac{p\,\big|\vv_{l,k}\herm\hat{\vh}_{l,k}\big|^{2}}
{p\sum_{j\neq k}\big|\vv_{l,k}\herm\hat{\vh}_{l,j}\big|^{2}+\vv_{l,k}\herm\mQ_{l}\vv_{l,k}}.
\label{eq:sinr}
\end{equation}
\vexp{$\gamma_{l,k}$ is the post-combining SINR of user $k$ at AP $l$; the numerator is the wanted power, the denominator is interference from the other users plus noise-and-error.}%
{Here $\gamma_{l,k}$ is how clean user $k$'s signal is after AP $l$'s filter: the top is the useful
signal power, the bottom is the leftover interference from the other $K\!-\!1$ users plus the noise
and estimation error. Unlike channel norm, $\gamma_{l,k}$ already knows about interference and
estimation error --- this is what makes P1 ``detector-aware''.}
Under large-scale fading decoding (LSFD) the central processor fuses the AP outputs, and the total
(all-AP) SINR of user $k$ is the sum of the per-AP SINRs \cite{nayebi2016lsfd}
\begin{equation}
\Gamma_{k}=\sum_{l=1}^{L}\gamma_{l,k},\qquad R_{k}=\log_{2}(1+\Gamma_{k}).
\label{eq:total}
\end{equation}
\vexp{$\Gamma_k$ is the total SINR over all $L$ APs; $R_k$ the achievable rate.}%
{Here $\Gamma_{k}$ is the best SINR user $k$ could get if \emph{all} APs served it, and $R_{k}$ is
the matching rate. Equation \eqref{eq:total} (SINRs add across APs under optimal LSFD fusion) is the
foundation of P1: clustering just decides which terms of this sum we keep.}

\section{Baseline rules (what we compare against)}
\textbf{CN (channel norm).} With a per-regime strength metric $\mu_{l,k}$
($\mu_{l,k}=\|\hat{\vh}^{\det}_{l,k}\|^{2}$ for near-field links \cite{cui2022channel}, $\mu_{l,k}=
\beta_{l,k}$ for far-field), keep the APs above a fraction of the strongest, plus the best AP
\cite{bjornson2020scalable}:
\begin{equation}
\Sk^{\mathrm{CN}}=\Big\{l:\mu_{l,k}\ge\tfrac{1}{\delta}\max_{l'}\mu_{l',k}\Big\}\cup\{\arg\max_{l}\mu_{l,k}\}.
\label{eq:cn}
\end{equation}
\vexp{$\mu_{l,k}$ is the channel-strength metric; $\beta_{l,k}$ the large-scale gain; $\delta>1$ the threshold divisor.}%
{Here $\mu_{l,k}$ measures how strong the channel is (its energy), $\beta_{l,k}$ is the large-scale
gain (path loss/shadowing), and $\delta>1$ sets the threshold: a larger $\delta$ keeps more APs. CN
only looks at strength, so it cannot tell a strong-but-interfered AP from a useful one.}

\noindent\textbf{IR (information rate, Mashdour--de~Lamare).} Score each link by an
interference-free rate and keep those above the average, plus the best AP \cite{mashdour2024clustering}:
\begin{equation}
\rho_{l,k}=\log_{2}\!\Big(1+\frac{p\,(\operatorname{tr}\mR_{l,k}-\operatorname{tr}\mC_{l,k})}{p\operatorname{tr}\mC_{l,k}+1}\Big),\quad
\Sk^{\mathrm{IR}}=\{l:\rho_{l,k}\ge c\,\bar\rho\}\cup\{\arg\max_{l}\rho_{l,k}\}.
\label{eq:ir}
\end{equation}
\vexp{$\rho_{l,k}$ is the per-link rate (signal over estimation error plus noise, no interference); $\mR_{l,k}$ the channel covariance; $\bar\rho$ the mean rate over all links; $c$ a scale factor.}%
{Here $\rho_{l,k}$ is a per-link rate built from the estimated-signal power $\operatorname{tr}\mR_{l,k}
-\operatorname{tr}\mC_{l,k}$ over the estimation error $\operatorname{tr}\mC_{l,k}$ plus noise; $\mR_{l,k}$
is the channel covariance, $\bar\rho$ the average of $\rho$ over all links, and $c$ a scaling knob.
\textbf{Limitation:} $\rho_{l,k}$ has \emph{no} interference term, and because it is an \emph{absolute}
threshold, more links pass as the SNR rises, so the cluster (and the fronthaul) grows with SNR.}

\section{Proposed rule: P1 (detector-aware information-surplus clustering)}
P1 works directly on the per-AP SINRs \eqref{eq:sinr}. Sort them in decreasing order,
\begin{equation}
\gamma_{(1),k}\ge\gamma_{(2),k}\ge\dots\ge\gamma_{(L),k},
\label{eq:sort}
\end{equation}
\vexp{$\gamma_{(i),k}$ is the $i$-th largest per-AP SINR of user $k$; $(\cdot)$ denotes the sorted order.}%
{Here $\gamma_{(i),k}$ is the $i$-th strongest AP for user $k$ (so $\gamma_{(1),k}$ is the best AP).
We rank the APs by how much clean signal each contributes.}
and keep the fewest top APs whose SINRs already cover a fraction $(1-\epsilon)$ of the total:
\begin{equation}
\boxed{\;m_{k}=\min\Big\{m:\ \textstyle\sum_{i=1}^{m}\gamma_{(i),k}\ \ge\ (1-\epsilon)\,\Gamma_{k}\Big\},
\qquad \Sk^{\mathrm{P1}}=\{(1),\dots,(m_{k})\}.\;}
\label{eq:p1}
\end{equation}
\vexp{$m_k=|\Sk^{\mathrm{P1}}|$ is the cluster size; $\epsilon\in(0,1)$ the surplus tolerance ($1-\epsilon$ the target coverage); $\Gamma_k$ from \eqref{eq:total}.}%
{Here $m_{k}$ is the chosen cluster size (number of APs for user $k$), and $\epsilon\in(0,1)$ is the
only knob: $1-\epsilon$ is the fraction of the total SINR we insist on keeping (e.g.\ $\epsilon=0.1$
keeps $90\%$). We add APs, strongest first, until their SINRs reach $(1-\epsilon)\Gamma_{k}$, then
stop. The best AP is always included, so $m_{k}\ge1$.}

\section{Theoretical guarantee}
Define the \emph{captured-information ratio} and the \emph{information surplus} (the normalized rate
the cluster leaves outside, eq.~(26) of \cite{haider2026vtc}):
\begin{equation}
\kappa_{k}=\frac{\Gamma_{\Sk}}{\Gamma_{k}},\qquad
\Delta_{k}=\frac{\Gamma_{k}-\Gamma_{\Sk}}{(1+\Gamma_{\Sk})(1+\Gamma_{k})},\qquad
\Gamma_{\Sk}=\sum_{l\in\Sk}\gamma_{l,k}.
\label{eq:kapdel}
\end{equation}
\vexp{$\kappa_k\in[0,1]$ is the fraction of SINR the cluster keeps; $\Delta_k\ge0$ the information lost outside the cluster; $\Gamma_{\Sk}$ the cluster SINR.}%
{Here $\kappa_{k}$ is the fraction of the total SINR the cluster keeps ($\kappa_{k}=1$ means nothing
is lost), $\Gamma_{\Sk}$ is the SINR summed over the serving APs only, and $\Delta_{k}$ is the
(small) amount of rate the user loses by not using every AP. P1 is designed to make $\kappa_{k}$
large and $\Delta_{k}$ small.}
By construction \eqref{eq:p1} gives $\Gamma_{\Sk}\ge(1-\epsilon)\Gamma_{k}$, i.e.\ $\kappa_{k}\ge
1-\epsilon$. This yields a per-user bound on the surplus:
\begin{propbox}
\textbf{Proposition (surplus bound).} For the P1 cluster \eqref{eq:p1},
\begin{equation}
\kappa_{k}\ge 1-\epsilon
\qquad\Longrightarrow\qquad
\Delta_{k}\ \le\ \frac{\epsilon}{\,1+(1-\epsilon)\Gamma_{k}\,}.
\label{eq:prop}
\end{equation}
\emph{Proof.} From $\Gamma_{\Sk}\ge(1-\epsilon)\Gamma_{k}$ we have $\Gamma_{k}-\Gamma_{\Sk}\le
\epsilon\Gamma_{k}$ and $1+\Gamma_{\Sk}\ge 1+(1-\epsilon)\Gamma_{k}$. Substituting in
\eqref{eq:kapdel} and using $1+\Gamma_{k}\ge\Gamma_{k}$,
\[
\Delta_{k}\le\frac{\epsilon\Gamma_{k}}{\big(1+(1-\epsilon)\Gamma_{k}\big)(1+\Gamma_{k})}
\le\frac{\epsilon\Gamma_{k}}{\big(1+(1-\epsilon)\Gamma_{k}\big)\Gamma_{k}}
=\frac{\epsilon}{1+(1-\epsilon)\Gamma_{k}}.\qquad\blacksquare
\]
\end{propbox}
\vexp{$\epsilon$ is the tolerance from \eqref{eq:p1}; $\Gamma_k$ the total SINR.}%
{\textbf{What this means in simple words:} the user never loses more than
$\epsilon/(1+(1-\epsilon)\Gamma_{k})$ of its rate by using the small cluster instead of all APs, and
this bound \emph{shrinks} as the user's SINR $\Gamma_{k}$ grows. So one knob $\epsilon$ directly
controls the worst-case information loss --- a guarantee that CN \eqref{eq:cn} and IR \eqref{eq:ir}
do not provide.}

\section{Cluster-size behaviour and comparison}
The size $m_{k}$ in \eqref{eq:p1} depends on how the SINR is spread over the APs: if one AP dominates
(far-field user with a single strong link) few APs reach $(1-\epsilon)\Gamma_{k}$, so $m_{k}$ is
small; if the useful SINR is spread over several nearby APs (near-field user) $m_{k}$ is larger.
Hence P1 \emph{adapts} the cluster to the user's geometry. Empirically (our study):
\begin{itemize}[leftmargin=1.4em,itemsep=1pt]
\item \textbf{Size:} $|\Sk^{\mathrm{CN}}|<|\Sk^{\mathrm{P1}}|<|\Sk^{\mathrm{IR}}|$ at high SNR
      ($\approx1.4$ vs $\approx2.0$ vs up to $3.6$).
\item \textbf{Stability:} $|\Sk^{\mathrm{P1}}|$ is \emph{flat in SNR} (it is a ratio target), whereas
      $|\Sk^{\mathrm{IR}}|$ \emph{grows with SNR} (absolute threshold) --- so IR's fronthaul grows
      while P1's stays bounded.
\item \textbf{User-type adaptivity:} P1 gives near-field users a slightly larger cluster than
      far-field users; CN treats both alike and IR over-serves near-field users.
\end{itemize}

\section{Benefits}
\begin{itemize}[leftmargin=1.4em,itemsep=1pt]
\item \textbf{Provable guarantee:} one knob $\epsilon$ bounds the per-user information loss
      \eqref{eq:prop}; neither CN nor IR offers this.
\item \textbf{Detector-aware:} P1 uses the post-combining SINR \eqref{eq:sinr}, which already contains
      interference and estimation error, unlike CN (strength only) and IR (no interference).
\item \textbf{Bounded, stable fronthaul:} the size target is a \emph{ratio}, so the cluster does not
      grow with SNR (unlike IR).
\item \textbf{Geometry-adaptive} cluster sizing (near-/far-field) and better AP load balance.
\item \textbf{Performance:} matches CN and clearly beats IR in fronthaul; on simple detectors
      (linear, SIC, list-SIC) it attains the lowest BER at high SNR. (On the cross-AP list detector
      all rules coincide, so the gain there is in fronthaul/interpretability, not BER.)
\end{itemize}

\section{Novelty and prior work}
Fraction/threshold AP selection on \emph{large-scale fading} is known: dynamic cooperation
clustering and largest-large-scale-fading selection \cite{bjornson2020scalable,ngo2018totalenergy},
user-centric clustering \cite{buzzi2017usercentric,interdonato2019ubiquitous}, and rate-threshold
clustering \cite{mashdour2024clustering}. \textbf{What is new in P1} is: (i) the selection criterion
is the \emph{post-combining per-AP SINR} \eqref{eq:sinr} (interference- and estimation-error aware),
not channel norm or an interference-free rate; and (ii) the cluster is sized by covering a fraction
$(1-\epsilon)$ of the \emph{total SINR} \eqref{eq:p1}, which yields the \emph{provable
information-surplus bound} \eqref{eq:prop} linking the single knob $\epsilon$ to the worst-case rate
loss. To our knowledge this SINR-coverage criterion with a surplus guarantee has not been proposed
for cell-free clustering; it should be positioned against \cite{bjornson2020scalable,ngo2018totalenergy,
mashdour2024clustering} and verified with a final literature check before claiming priority.

\begin{keybox}
\textbf{Honest scope.} The information-loss guarantee \eqref{eq:prop} is exact under the
SINRs-add model \eqref{eq:total}. End-to-end, P1 matches channel-norm clustering on the proposed
cross-AP list detector (the detector absorbs the clustering choice) and its advantage there is
\emph{bounded fronthaul and interpretability}; its clearest gains are the surplus guarantee, the
fronthaul saving over IR, and lower BER on the simpler detectors. The clustering step costs more
computation than CN (it needs the per-AP SINRs), which should be reported.
\end{keybox}

\begin{thebibliography}{99}\small
\bibitem{ngo2017cellfree} H.~Q.~Ngo \emph{et al.}, ``Cell-free massive MIMO versus small cells,'' \emph{IEEE Trans. Wireless Commun.}, 16(3):1834--1850, 2017.
\bibitem{ngo2018totalenergy} H.~Q.~Ngo \emph{et al.}, ``On the total energy efficiency of cell-free massive MIMO,'' \emph{IEEE Trans. Green Commun. Netw.}, 2(1):25--39, 2018.
\bibitem{bjornson2020scalable} E.~Bj\"ornson and L.~Sanguinetti, ``Scalable cell-free massive MIMO systems,'' \emph{IEEE Trans. Commun.}, 68(7):4247--4261, 2020.
\bibitem{interdonato2019ubiquitous} G.~Interdonato \emph{et al.}, ``Ubiquitous cell-free massive MIMO communications,'' \emph{EURASIP J. Wireless Commun. Netw.}, 2019:197, 2019.
\bibitem{buzzi2017usercentric} S.~Buzzi and C.~D'Andrea, ``Cell-free massive MIMO: User-centric approach,'' \emph{IEEE Wireless Commun. Lett.}, 6(6):706--709, 2017.
\bibitem{nayebi2016lsfd} E.~Nayebi \emph{et al.}, ``Performance of cell-free massive MIMO systems with MMSE and LSFD receivers,'' in \emph{Proc. Asilomar}, 2016, pp.~203--207.
\bibitem{mashdour2024clustering} S.~Mashdour, R.~C.~de~Lamare, and J.~P.~S.~H.~Lima, ``Clustering and resource allocation for cell-free massive MIMO systems,'' arXiv:2404.18032, 2024 (verify exact title/venue).
\bibitem{cui2022channel} M.~Cui and L.~Dai, ``Channel estimation for extremely large-scale MIMO: Far-field or near-field?'' \emph{IEEE Trans. Commun.}, 70(4):2663--2677, 2022.
\bibitem{haider2026icassp} N.~Haider and R.~C.~de~Lamare, ``Cross-AP list-based interference cancellation for near-field cell-free XL-MIMO,'' submitted to IEEE ICASSP, 2026.
\bibitem{haider2026vtc} N.~Haider and R.~C.~de~Lamare, ``Cross-AP channel estimation and interference cancellation for near-field cell-free XL-MIMO,'' submitted to IEEE VTC, 2026.
\bibitem{haider2026wcl} N.~Haider and R.~C.~de~Lamare, ``Iterative interference cancellation for mixed near- and far-field XL-MIMO,'' \emph{IEEE WCL}, 15:5323--5327, 2026.
\end{thebibliography}
\end{document}
